\section{The Navigation-Primitive Taxonomy}
\label{sec:taxonomy}

Table~\ref{tab:comparison} shows the five primitives. In brief:
\src{P} walks a filesystem tree (name~$\to$~inode~$\to$~block);
\src{M} walks a kernel object graph (PID~$\to$~EPROCESS~$\to$~VA~$\to$~PA);
\src{L} seeks a record stream by timestamp or cursor;
\src{Q} dispatches a query to an endpoint and receives result rows that did
not exist before the query was issued;
\src{C} follows cryptographic hashes through a Merkle DAG where the address
\emph{is} the hash of the content.
The formal definitions below make these distinctions precise.

\begin{table*}[!t]
\centering
\caption{The Five Navigation Primitives}
\label{tab:comparison}
\begin{tabular}{@{}p{0.06\textwidth}p{0.18\textwidth}p{0.20\textwidth}p{0.20\textwidth}p{0.10\textwidth}p{0.14\textwidth}@{}}
\toprule
\textbf{Type} & \textbf{Address space} & \textbf{Traversal operation} & \textbf{Canonical tools} & \textbf{Mutability} & \textbf{Attacker durability} \\
\midrule
\src{P} & path / inode / block & directory walk; inode lookup; block read & TSK, Autopsy, ewf, ext4fs-forensic & mutable on host & by acquisition \\
\src{M} & (PID, VA) $\to$ PA & process-list walk; page-table walk & Volatility, Rekall, memf-* & mutable on host (DKOM) & by acquisition \\
\src{L} & timestamp / cursor & seek to record boundary; field decode & Plaso, EvtxECmd, winevt-forensic & append-only by intent & by acquisition \\
\src{Q} & (endpoint, query, cursor) & query dispatch; cursor pagination & Velociraptor, OSQuery, WMI & N/A (produced, not stored) & by transmission \\
\src{C} & cryptographic hash & hash $\to$ blob; reference following & git, Sigstore, OCI, IPFS & immutable by construction & intrinsic \\
\bottomrule
\end{tabular}
\end{table*}

\subsection{Formal Foundations}

We treat byte sequences, denoted $\Bstar = \{0,1,\ldots,255\}^{*}$, as the
unstructured medium through which all forensic evidence ultimately passes.
Forensic records, denoted $R$, are the typed objects that populate analyst
reports, timelines, and ontologies (\eg{}, browser history rows, EVTX event
records, EPROCESS structures). Parsers translate from $\Bstar$ to $R$.
Navigation translates from addresses to $\Bstar$. We work in the category
$\mathbf{PFn}$ of sets and partial functions~\cite{maclane1998categories}:
objects are sets, morphisms are partial functions, composition is the usual
chain composition with the convention that an undefined intermediate value
propagates as undefined.

\begin{definition}[Evidence Source]
\label{def:source}
An \emph{evidence source} is a triple $S = (A,\, f_{\mathrm{nav}},\,
f_{\mathrm{parse}})$, where:
\begin{itemize}
    \item $A$ is an \emph{address space}: a set whose elements name reachable
    locations within the source;
    \item $f_{\mathrm{nav}} : A \rightharpoonup \Bstar$ is a (partial)
    \emph{navigation function} that, given an address, yields the byte
    sequence stored at that address; and
    \item $f_{\mathrm{parse}} : \Bstar \to R$ is a \emph{parser} that
    interprets bytes as forensic records.
\end{itemize}
The composition $f_{\mathrm{parse}} \circ f_{\mathrm{nav}} : A \rightharpoonup R$
yields records keyed by source address.
\end{definition}

\begin{definition}[Address-Space Algebra]
\label{def:addr-algebra}
An \emph{address-space algebra} is a tuple $\mathcal{A} = (A, \mathrm{Op}, \to)$
where $A$ is a set of addresses, $\mathrm{Op}$ is a finite signature of
\emph{traversal operations} (each $\omega \in \mathrm{Op}$ is a partial
function $\omega : A^{k_\omega} \to A$ for some arity $k_\omega \geq 0$),
and $\to \;\subseteq A \times A$ is the reachability relation generated by
$\mathrm{Op}$ (\ie{}, $a \to a'$ iff some $\omega \in \mathrm{Op}$ maps $a$
to $a'$, and $\to$ is taken under transitive-reflexive closure).
\end{definition}

\begin{definition}[Navigation Primitive]
\label{def:nav-primitive}
A \emph{navigation primitive} $\pi$ is an isomorphism class of address-space
algebras under the equivalence relation $\sim$ defined as follows. Two
algebras $\mathcal{A} = (A, \mathrm{Op}, \to)$ and $\mathcal{A}' = (A',
\mathrm{Op}', \to')$ are equivalent, $\mathcal{A} \sim \mathcal{A}'$, iff
there exist a bijection $\varphi : A \to A'$ and a one-to-one
correspondence $\sigma : \mathrm{Op} \to \mathrm{Op}'$ that preserves arity
and that satisfies $\varphi(\omega(a_1, \ldots, a_{k_\omega})) =
\sigma(\omega)(\varphi(a_1), \ldots, \varphi(a_{k_\omega}))$ wherever the
left-hand side is defined. The bijection condition guarantees no
information loss; the operation correspondence guarantees that the same
traversal questions can be answered in either algebra.
\end{definition}

\begin{definition}[Source Type]
\label{def:source-type}
Two evidence sources $S = (A, f_{\mathrm{nav}}, f_{\mathrm{parse}})$ and
$S' = (A', f'_{\mathrm{nav}}, f'_{\mathrm{parse}})$ are of the same
\emph{source type} iff their underlying address-space algebras are equivalent
under $\sim$ (Definition~\ref{def:nav-primitive}), \ie{}, they share a
navigation primitive $\pi$.
\end{definition}

Two sources are the same type iff their addressing question -- ``starting at
$a$, where can I go and by what operation?'' -- is substitutable. An EVTX
file and a syslog stream are both \src{L}. A Chrome \texttt{History} file on
a mounted disk (\src{P}) and the same file recovered from carved memory (\src{M})
are not the same type, despite being the same artifact.

\subsection{The Five Primitives}

\subsubsection{\src{P} -- Persistent Storage}

\paragraph{Address space.} The filesystem namespace, $A_{P} = \mathrm{Path}
\cup \mathrm{Inode} \cup \mathrm{Block}$, with structure-preserving maps
$\mathrm{Path} \to \mathrm{Inode} \to \mathrm{Block}^{*}$ supplied by the
filesystem.
\paragraph{Traversal.} Resolve a path to an inode by directory walk; resolve
the inode to a block list by inode-table lookup; resolve the block list to
bytes by issuing reads against the underlying container (sector stream from
an E01 image, raw partition, FUSE mount, \etc{}).
\paragraph{Canonical examples.} E01/EWF disk images~\cite{ewf2011}, raw
partitions, NTFS, ext4, APFS, FUSE-bridged images.
\paragraph{Epistemological properties.} High completeness when the volume is
acquired offline; vulnerable to anti-forensic deletion and timestomping;
records are mutable in principle but typically frozen by acquisition.

\subsubsection{\src{M} -- Memory}

\paragraph{Address space.} Per-process virtual address space $A_{M} =
\mathrm{PID} \times \mathrm{VA}$, layered over the OS object graph and
ultimately resolving through page tables to physical addresses
$\mathrm{PA}$.
\paragraph{Traversal.} Walk the kernel's process list to reach an
\texttt{EPROCESS} (Windows) or \texttt{task\_struct} (Linux); walk the VAD
tree (Windows) or VMA list (Linux) to enumerate mapped regions; walk page
tables (PML4 on x86\_64, PAE on legacy x86, the four-level scheme on
AArch64) to map a virtual address to a physical
page~\cite{ligh2014artmemory,case2017memory}.
\paragraph{Canonical examples.} WinPMEM, LiME, AVML, hiberfil.sys,
Windows crash dumps.
\paragraph{Epistemological properties.} Captures non-persistent state
unavailable from \src{P}; vulnerable to direct kernel object manipulation
(DKOM); contents are an instantaneous snapshot rather than a trace.

\subsubsection{\src{L} -- Log}

\paragraph{Address space.} A linearly ordered space of cursors, $A_{L}
\subseteq \mathbb{N} \cup \mathrm{Timestamp} \cup \mathrm{Cursor}$, where a
cursor names a record boundary within a stream.
\paragraph{Traversal.} Seek to a chunk or block boundary by timestamp or
sequence number; read forward, decoding record boundaries; decode each record
to a structured field set. Backward traversal requires either an index or
linear scan.
\paragraph{Canonical examples.} Windows EVTX binary
format~\cite{microsoft2023evtx}, systemd-journald binary
journal~\cite{poettering2012journal}, Apple Unified Logging
\texttt{tracev3}~\cite{apple2016unifiedlogging}, AWS CloudTrail JSON event
streams~\cite{aws2024cloudtrail}, PCAP-NG capture
files~\cite{pcapng2024}.
\paragraph{Epistemological properties.} Append-only by intent (mutation
without trace requires privileged on-host action); linearly ordered, which
greatly simplifies temporal correlation; subject to rollover and rotation
that can lose records.

\subsubsection{\src{Q} -- Live Query}

\paragraph{Address space.} The Cartesian product of well-formed query
expressions and pagination cursors, $A_{Q} = \mathrm{Endpoint} \times
\mathrm{Query} \times \mathrm{Cursor}$.
\paragraph{Traversal.} Dispatch a query to an endpoint; the endpoint
executes the query against live state; results stream back as typed rows,
with cursor tokens advancing pagination. Re-execution of the same query at a
different time generally yields a different result.
\paragraph{Canonical examples.} Velociraptor VQL~\cite{velociraptor},
WMI/WQL~\cite{microsoft2023wmi}, OSQuery SQL~\cite{osquery}, GRR
Rapid Response~\cite{cohen2011distributed}.
\paragraph{Epistemological properties.} Evidence is \emph{produced} by the
act of querying; it does not pre-exist in any addressable store. Once
captured to an investigator-controlled system, the result is durable against
on-host adversarial action (Section~\ref{sec:properties}). The query itself
is part of the evidence chain and must be recorded.

\subsubsection{\src{C} -- Content-Addressed}

\paragraph{Address space.} The image of a cryptographic hash function, $A_{C}
\subseteq \{0,1\}^{n}$ for hash output width $n$ (\eg{}, SHA-1, SHA-256). The
\emph{address is the hash}.
\paragraph{Traversal.} Given a hash, retrieve the blob whose hash is that
value; parse the blob's references (each itself a hash); follow references
to traverse a Merkle DAG. Cycles are impossible because each node's hash
depends on the hashes of its references.
\paragraph{Canonical examples.} Git commit/tree/blob graphs~\cite{chacon2014progit},
OCI image manifests and layers~\cite{ociimagespec}, IPFS~\cite{benet2014ipfs},
the Sigstore Rekor transparency log~\cite{newman2022sigstore},
in-toto links~\cite{torresarias2019intoto}.
\paragraph{Epistemological properties.} Identity equals content. Modification
yields a new address rather than altering the artifact at the original
address (Section~\ref{sec:properties}). Enables \emph{hash-pivot}:
simultaneous query of all known CAS stores by a single hash.

\subsection{Layer Hierarchy}

The taxonomy yields a five-layer architecture in which each layer has a
defined dependency relationship with the layers below. Figure~\ref{fig:layers}
sketches the dependencies; Section~\ref{sec:implementation} maps these
layers to concrete crates in the Issen reference implementation.

\begin{figure}[t]
\centering
\begin{tabular}{@{}l@{}}
\toprule
\textbf{ORCHESTRATION}~~ wires sources, correlates evidence, drives UX \\
~~$\uparrow$ depends on all layers below \\
\midrule
\textbf{PARSER}~~~~~~~~~~~ $\Bstar \to R$, source-agnostic \\
~~$\uparrow$ depends only on KNOWLEDGE \\
\midrule
\textbf{NAVIGATE}~~~~~~~~~ $A \to \Bstar$, primitive-specific\\
~~ \src{P}: filesystem ~~~~~ \src{M}: paging / OS structure \\
~~ \src{L}: log format ~~~~~~ \src{Q}: query engine \\
~~ \src{C}: graph navigation \\
~~$\uparrow$ depends on CONTAINER + KNOWLEDGE \\
\midrule
\textbf{CONTAINER}~~~~~~ raw stream from outer format (E01, dump, \ldots) \\
~~$\uparrow$ depends on KNOWLEDGE \\
\midrule
\textbf{KNOWLEDGE}~~~~ format constants, magics, schemas \\
\bottomrule
\end{tabular}
\caption{Five-layer architecture induced by the taxonomy. Higher layers
depend only on lower layers. PARSER is independent of NAVIGATE because both
operate at the byte-sequence interface (Proposition~\ref{thm:parser-indep}).}
\label{fig:layers}
\end{figure}

The KNOWLEDGE layer holds purely declarative facts about formats: page
sizes, magic numbers, struct field offsets, schema definitions. The
CONTAINER layer decodes outer dump or image formats (E01 segments, WinPMEM
output, ELF core dumps) into a raw addressable stream; \src{Q} and \src{C}
have no separate container layer because the endpoint or hash store \emph{is}
the entry point. The NAVIGATE layer implements a navigation primitive
against a stream or store. The PARSER layer interprets bytes as records and
is, by construction, independent of which navigation primitive produced
those bytes. The ORCHESTRATION layer is the only place that may depend on
all of the above, and is responsible for routing artifacts to parsers and
for cross-source correlation.

We will see in Section~\ref{sec:properties} that the dependency rules are not
merely a stylistic convention; they are forced by Proposition~\ref{thm:parser-indep}
and by the contrapositive of Proposition~\ref{thm:layer-rule}.
